\title{Direct Preference Optimization: Your Language Model is Secretly a Reward Model}

\begin{abstract}
Although large-scale unsupervised language models (LMs) acquire extensive world knowledge and display some reasoning capabilities, precisely directing their outputs remains challenging due to the fully unsupervised framework of their training processes. To address this limitation, prevalent approaches employ human-annotated data indicating relative preferences between model outputs, subsequently fine-tuning the unsupervised LM in accordance with these preferences, often through reinforcement learning from human feedback (RLHF). Nevertheless, RLHF is typically a multi-stage and potentially unstable approach that first constructs a reward model representing human preferences, and then applies reinforcement learning to the LM so as to optimize this modeled reward without diverging excessively from the initial model behavior. In this work, we propose a novel way to parameterize the reward model used in RLHF, which makes it possible to analytically determine the associated optimal policy, thereby reducing the RLHF task to optimizing a straightforward classification objective. The resulting approach, termed \textit{Direct Preference Optimization} (DPO), is robust, effective, and computationally efficient, as it dispenses with the need to generate samples from the LM during fine-tuning or perform extensive hyperparameter searches. Through a set of empirical studies, we demonstrate that DPO aligns LMs with human preferences as well as, or better than, prior techniques. Remarkably, DPO-based fine-tuning surpasses RLHF methods employing PPO in guiding generation sentiment, and achieves comparable or superior results in tasks such as summarization and single-turn dialogue, all while being significantly easier to implement and train.
\end{abstract}

\section{Introduction}
Large unsupervised language models (LMs) trained on massive corpora demonstrate a range of unexpected competencies~\citep{chowdhery2022palm, brown2020language, touvron2023llama, bubeck2023sparks}. Nonetheless, the training data for these models consists of human-generated text reflecting diverse motivations, intentions, and expertise levels. Not all such motivations or skill expressions are desirable for imitation; for instance, it is beneficial for an AI programming assistant to \textit{recognize} common coding errors to enable correction, but when producing code, we prefer the model to emulate the (sometimes infrequent) exemplary programming showcased in its training set. Analogously, although a language model should be \textit{informed} about widely held misconceptions (e.g., believed by half the population), it should not reproduce these inaccuracies proportionally in its responses. Thus, it is essential to \emph{select} the model’s \emph{preferred behaviors and answers} from its broad \textit{repertoire of knowledge and skills} to achieve AI systems that are reliable, effective, and amenable to user control \citep{ouyang2022training}. Existing strategies for steering language models toward preferred outputs frequently rely on reinforcement learning (RL) guided by human feedback, but as we discuss, the RL-based objective used in such approaches can, in fact, be solved directly using a binary cross-entropy loss, providing a much simpler route for preference alignment.

Generally, prevailing approaches guide language models to exhibit preferred behaviors by leveraging collections of human-labeled preferences characterizing what users find safe and beneficial. This stage of learning from preferences follows extensive unsupervised pre-training on voluminous textual data. Although supervised fine-tuning with high-quality human demonstrations is the most direct method for imparting preferences, the approaches that achieve the strongest empirical results are based on reinforcement learning from human or AI feedback (RLHF/RLAIF; \citep{christiano2017deep, bai2022constitutional}). In RLHF, a reward function is constructed by training a model on datasets of human judgments, and the language model is subsequently optimized via RL to increase reward, all the while avoiding excessive divergence from the base model. Despite the impressive results achieved by RLHF for both dialog and code generation tasks, the methodology remains significantly more elaborate than standard supervised fine-tuning. It requires the sequential training of multiple models and repeated sampling from the evolving model’s policy during optimization, leading to substantial computational demands.

This work introduces a novel approach for aligning language models with human preferences without the need for explicit reward models or reinforcement learning frameworks. We present \textit{Direct Preference Optimization} (DPO), an algorithm that inherently targets the same optimization goal as standard RLHF approaches—reward maximization with a KL-divergence regularization—while offering greater simplicity in both implementation and training procedures. DPO intuitively modifies the update rule to boost the log-likelihood ratio between favored and disfavored responses. Crucially, this process incorporates an adaptive, per-example weighting mechanism that addresses the model deterioration issues observed with a straightforward probability ratio objective. Similar to established techniques, DPO leverages a formal preference model (such as the Bradley-Terry model; \cite{bradley1952rankanalysis}) to quantify alignment between reward functions and observed preference data. However, unlike traditional pipelines—which use this model to derive a preference loss for training a reward model and subsequently update the policy to maximize rewards—DPO reframes the preference loss directly as a function of the policy using a change of variables. Consequently, given a corpus of preference-labeled examples, DPO enables direct policy optimization via a binary cross entropy loss, resulting in an optimal policy under an implicit reward inferred from the observed preferences.

The principal contribution of this work is the introduction of Direct Preference Optimization (DPO), a streamlined, reinforcement-learning-free technique for preference-based training of language models. Empirical results demonstrate that DPO matches or surpasses the effectiveness of prior methods—including PPO-based RLHF—across preference-driven tasks like sentiment adjustment, summarization, and conversational modeling, using models with parameters scaling up to 6B.

\section{Related Work}

Progressively larger self-supervised language models have demonstrated the ability to perform certain tasks with no additional training \citep{radford2019language}, or by leveraging a handful of in-context examples through few-shot prompting \citep{gpt3,megatron,chowdhery2022palm}. Yet, their effectiveness on specific downstream tasks and their capacity to follow user intent are notably enhanced when these models undergo fine-tuning using datasets composed of instructions paired with human-generated responses \citep{mishra-etal-2022-cross,sanh2022multitask,chung2022scaling,thoppilan2022lamda}. This process, known as ‘instruction-tuning,’ facilitates the extension of LLMs’ capabilities to encompass instructions that are not present in the tuning set, broadly improving their functionality and generalizability \citep{chung2022scaling}. Although instruction tuning has achieved considerable advances, collecting \textit{relative} human preferences between model outputs tends to be less resource-intensive than assembling high-quality expert demonstrations. Consequently, recent efforts have adopted preference-based fine-tuning, where LLMs are adapted using datasets annotated with human preferences. This strategy has demonstrated gains across various applications including translation \citep{kreutzer-etal-2018-reliability}, summarization \citep{stiennon2022learning,ziegler2020finetuning}, narrative generation \citep{ziegler2020finetuning}, and instruction following \citep{ouyang2022training,ramamurthy2023is}. Such methods proceed by first training a neural reward model to fit the preference-labeled data, often leveraging preference models such as the Bradley-Terry framework \citep{bradley1952rankanalysis}, and subsequently adjusting the language model itself to optimize this learned reward signal. Optimization is typically performed via reinforcement learning techniques including REINFORCE \citep{williams1992reinforce}, proximal policy optimization (PPO; \cite{schulman2017proximal}), or related approaches \citep{ramamurthy2023is}. Another related approach employs instruction-following LLMs, further refined through human feedback, as generators of synthetic preference data for desired characteristics such as safety or non-toxicity \citep{bai2022constitutional}. In these cases, models are weakly supervised with rubric-like guidance supplied by humans to direct the LLMs’ evaluation of outputs. Altogether, these techniques synthesize advances from two distinct but converging domains: reinforcement learning-based language model training for various tasks~\citep{Ranzato2015SequenceLT,paulus2018a,wu2018learning}, and general methods for preference learning from human feedback \citep{christiano2017deep,kupcsik2018learning}. Nevertheless, despite the attractiveness of learning from relative human preferences, applying reinforcement learning to fine-tune large language models introduces substantial operational complexities. The present work offers an approach grounded in theoretical principles for optimizing language models based on relative preferences, while avoiding the use of reinforcement learning.

Learning policies from preferences, outside of linguistic applications, has been a subject of investigation in both bandit and reinforcement learning paradigms, resulting in the development of diverse methodologies. When actions are ranked or compared via preferences instead of receiving scalar rewards in a contextual bandit framework, the problem is referred to as the contextual dueling bandit (CDB; \cite{yue2012karmed, dudik2015contextual}). Due to the lack of absolute rewards, theoretical work on CDBs adopts the concept of a \textit{von Neumann winner}—a policy that, in expectation, wins at least half the time against any alternative policy \citep{dudik2015contextual}. A crucial distinction is that CDBs usually acquire preference data in an online fashion, whereas human preference learning typically relies on a predetermined, finite dataset of actions paired with preference annotations collected offline \citep{yan2022human}. Similarly, in \textit{preference-based RL} (PbRL), the agent learns from binary preferences that reflect an unknown ‘scoring’ function, instead of rewards \citep{BusaFekete2014, ruiz2023dueling}. There is a range of PbRL algorithms, including those leveraging off-policy preference data, yet the majority explicitly estimate the underlying reward or scoring function before using it for policy optimization \citep{jain2013learning, BusaFekete2014, christiano2017deep, sadigh2017active, kupcsik2018learning}. Contrasting with these approaches, we introduce a one-stage method that aims to train policies directly to align with observed preferences, bypassing reward model estimation.

\section{Preliminaries}\label{section:prelims}

We summarize the RLHF framework described by \citeauthor{ziegler2020finetuning} and followed in later works (\citep{stiennon2022learning, bai2022training, ouyang2022training}). This process generally involves three steps: (1) supervised fine-tuning (SFT); (2) collecting preferences and fitting a reward model; and (3) reinforcement learning-based policy refinement.

\textbf{SFT}: The initial phase of RLHF employs supervised learning to further tune a large pre-trained language model on carefully curated datasets specific to downstream applications (e.g., dialogue, summarization), resulting in a policy $\pi^\text{SFT}$.

\textbf{Reward Modelling Phase}: Next, prompts $x$ are used to generate candidate completions $(y_1, y_2)\sim \pi^\text{SFT}(y \mid x)$ via the SFT model. Human annotators are then tasked with selecting a preferred response from these candidates, recorded as $y_w\succ y_l \mid x$ where $y_w$ denotes the favored and $y_l$ the less favored output. These judgments are presumed to derive from an underlying but unknown reward mechanism $r^*(y, x)$, which remains hidden. Among various statistical preference models, the Bradley-Terry (BT) model \cite{bradley1952rankanalysis} is widely employed (though more general frameworks such as Plackett-Luce models \citep{plackett1975analysis, luce2012individual} can also be applied when multiple outputs are ranked). In the BT framework, human preferences $p^*$ are parameterized as:

Given a fixed dataset of comparisons $\mathcal{D} = \bigl\{x^{(i)}, y_w^{(i)}, y_l^{(i)}\bigr\}_{i=1}^N$ drawn from $p^*$, one can introduce a reward function $r_{\phi}(x, y)$ parameterized by $\phi$, and estimate these parameters via maximum likelihood estimation. This task can be cast as a binary classification problem, leading to the following negative log-likelihood objective:

\begin{equation}\label{eq:reward_model}
    \mathcal{L}_R(r_{\phi}, \mathcal{D}) = -\mathbb{E}_{(x, y_w, y_l)\sim \mathcal{D}}\bigl[\log \sigma(r_{\phi}(x, y_w) - r_{\phi}(x, y_l))\bigr]
\end{equation}

where $\sigma$ denotes the logistic sigmoid function. For language models, $r_{\phi}(x, y)$ is typically initialized using the SFT model $\pi^\text{SFT}(y \mid x)$, with an additional linear head atop the final transformer representation to output a scalar reward value~\cite{ziegler2020finetuning}. To control reward variance, previous work applies normalization such that $\mathbb{E}_{x,y\sim \mathcal{D}}\left[r_\phi(x, y)\right]=0$ over the data.

\textbf{RL Fine-Tuning Phase}: Once the reward model is learned, it serves as a source of signal during fine-tuning. The optimization objective, as established in earlier works~\citep{jaques2017sequence, jaques2020human}, is given by

\begin{equation}\label{eq:RL}
\max_{\pi_{\theta}}  \mathbb{E}_{x\sim \mathcal{D}, y\sim \pi_{\theta}(y \mid x)}\bigl[r_{\phi}(x, y)\bigr] - \beta\mathbb{D}_{\textrm{KL}}\bigl[\pi_{\theta}(y\mid x)\mid \mid \pi_\text{ref}(y\mid x)\bigr],
\end{equation}

with $\beta$ modulating the penalty for deviating from the reference distribution $\pi_\text{ref}$, which is usually instantiated as the SFT policy $\pi^\text{SFT}$. Conventionally, the policy $\pi_\theta$ is also initialized from $\pi^\text{SFT}$. Incorporating this KL regularization is crucial; it limits policy drift beyond the region where the reward model is dependable and supports response diversity, thereby mitigating mode collapse towards a few high-reward sequences. Due to the inherently discrete generation process in language modeling, the above objective is non-differentiable and generally requires reinforcement learning methods for optimization. The commonly adopted strategy~\citep{ziegler2020finetuning, stiennon2022learning, bai2022training, ouyang2022training} involves forming the composite reward $r(x, y) = r_{\phi}(x, y) - \beta (\log \pi_{\theta}(y\mid x) - \log \pi_\text{ref}(y\mid x))$ and optimizing it via PPO~\cite{schulman2017proximal}.

\section{Direct Preference Optimization}\label{sec:DPO}

Recognizing the difficulties posed by reinforcement learning methods when scaling to large problems such as language model fine-tuning, we set out to develop a method for direct policy optimization using preference data. Departing from traditional RLHF pipelines, which fit an explicit reward function and subsequently optimize it, our approach exploits a specialized parameterization of the reward model. This allows for analytical recovery of the corresponding optimal policy without recourse to an iterative RL procedure. Central to our methodology is an explicit, closed-form relationship between reward models and their induced optimal policies. This enables conversion of an optimization objective defined on reward functions into one defined directly on policies, utilizing a change-of-variables. The resulting approach obviates the need for a separately-trained reward model, while continuing to faithfully represent established frameworks for human preference modeling, such as the Bradley-Terry model. In this formulation, the policy model encapsulates both the language generation mechanism and the implicit reward signal.

\textbf{Deriving the DPO objective.} As in previous studies, we begin with the RL objective provided in Eq.~\ref{eq:RL}, assuming an arbitrary reward function $r$. Building on the results of previous literature~\citep{peters2007reinforcement, peng2019advantage, korbak2022reinforcement, go2023aligning}, one can readily derive that the optimal policy for the KL-regularized reward maximization problem in Eq.~\ref{eq:RL} has the following structure:

\begin{equation}\label{eq:op_policy}
    \pi_r(y\mid x) = \frac{1}{Z(x)}\pi_\text{ref}(y\mid x)\exp\left(\frac{1}{\beta}r(x, y)\right),
\end{equation}

Here, $Z(x) =\sum_{y}\pi_\text{ref}(y\mid x)\exp\left(\frac{1}{\beta}r(x, y)\right)$ serves as the normalization constant (partition function). Details on this derivation can be found in Appendix~\ref{app:derivation1}. Estimating $Z(x)$ remains computationally demanding even if the reward function $r^*$ is replaced with its MLE approximation $r_{\phi}$, which poses practical challenges~\citep{korbak2022reinforcement, go2023aligning}. Nevertheless, by manipulating Eq.~\ref{eq:op_policy}, the reward can be re-expressed as a function of the optimal policy $\pi_r$, the baseline $\pi_\text{ref}$, and the partition function $Z(\cdot)$. Explicitly, by applying the logarithm to both sides of Eq.~\ref{eq:op_policy}, followed by some algebraic rearrangement, we derive:

\begin{equation}\label{eq:main_eq}
    r(x,y) =\beta \log \frac{\pi_r(y\mid x)}{\pi_\text{ref}(y\mid x)} + \beta \log Z(x).
\end{equation}

This change of variables may then be applied for the true reward $r^*$ and its corresponding optimal policy $\pi^*$. Importantly, the Bradley-Terry preference model relies exclusively on the difference between rewards for two responses, namely, ${p^*(y_1 \succ y_2 \mid x) = \sigma(r^*(x, y_1) - r^*(x, y_2))}$. Plugging the parameterization from Eq.~\ref{eq:main_eq} into the preference expression Eq.~\ref{eq:bradley-terry} leads to cancellation of the partition function. This enables the preference probability to be written solely in terms of the optimal policy $\pi^*$ and the reference policy $\pi_\text{ref}$. Consequently, the optimal RLHF policy $\pi^*$ under the Bradley-Terry model must satisfy:

\begin{equation}\label{eq:objective}
    p^*(y_1\succ y_2 \mid x)=\frac{1}{1 + \exp\left(\beta \log \frac{\pi^*(y_2\mid x)}{\pi_\text{ref}(y_2\mid x)} - \beta \log \frac{\pi^*(y_1\mid x)}{\pi_\text{ref}(y_1\mid x)}\right)}
\end{equation}

A thorough derivation is available in Appendix~\ref{app:derivation2}. Although Eq.~\ref{eq:objective} centers on the Bradley-Terry approach, an analogous methodology can be utilized for deriving preference models based on the more general Plackett-Luce models~\citep{plackett1975analysis, luce2012individual}; refer to Appendix~\ref{app:plackett_luce_models} for details.

With this expression for the probability of human preferences in terms of the optimal policy, rather than the reward, we are able to establish a maximum likelihood training objective for any parametric policy $\pi_\theta$. In direct analogy to standard reward modeling (cf. Eq.~\ref{eq:reward_model}), our new objective for the policy becomes:

\begin{equation}\label{eq:optimum_model}
    \mathcal{L}_\text{DPO}(\pi_{\theta}; \pi_\text{ref}) = -\mathbb{E}_{(x, y_w, y_l)\sim \mathcal{D}}\left[\log \sigma \left(\beta \log \frac{\pi_{\theta}(y_w\mid x)}{\pi_\text{ref}(y_w\

In this approach, we model an implicit reward through an alternative parametrization, wherein the optimal policy is directly represented by $\pi_\theta$. Additionally, because our method corresponds to a reparameterized Bradley-Terry model, it inherits specific theoretical guarantees—such as consistency, provided certain assumptions about the distribution of the preference data are met \cite{bong2022generalized}. Further theoretical considerations regarding DPO, including its connection to related methods, are provided in Section~\ref{sec:theory}.

\textbf{What is the effect of the DPO update?} To better grasp the mechanism underlying DPO, we can study the gradient of the loss function $\mathcal{L}_\text{DPO}$. Specifically, the gradient with respect to $\theta$ takes the following form:

\begin{multline*}\label{eq:gradient}
    \nabla_\theta \mathcal{L}_\text{DPO}(\pi_\theta;\pi_\text{ref}) = \\ -\beta\mathbb{E}_{(x, y_w, y_l) \sim \mathcal{D}} \bigg[\underbrace{\sigma(\hat{r}_\theta(x, y_l) - \hat{r}_\theta (x, y_w))}_\text{larger for incorrectly estimated rewards}\bigg[\underbrace{\nabla_\theta\log \pi(y_w \mid x)}_\text{boost probability of $y_w$} - \underbrace{\nabla_\theta\log\pi(y_l \mid x)}_\text{reduce probability of $y_l$}\bigg]\bigg],
\end{multline*}

where $\hat{r}_\theta(x, y) = \beta \log \frac{\pi_\theta(y \mid x)}{\pi_\text{ref}(y \mid x)}$ denotes the reward implicitly specified by the language model $\pi_\theta$ relative to the reference model $\pi_\text{ref}$ (for additional discussion, see Section~\ref{sec:theory}). From an intuitive perspective, this gradient promotes increasing the likelihood of the preferred response $y_w$ while simultaneously lowering that of the dispreferred response $y_l$. The influence of each example is modulated by the extent to which the implicit reward function $\hat{r}_\theta$ overvalues the dispreferred response $y_l$ compared to $y_w$, with $\beta$ adjusting for the effect of the KL regularization. This weighting reflects the model’s error in ranking the responses. Our empirical findings underscore the criticality of this weighted scheme, as omitting the weighting term can result in collapse or degraded language model quality (see Appendix Table~\ref{tab:unlikelihood_generations}).

\textbf{DPO framework.} The DPO procedure typically involves two main steps: 1) For each input prompt $x$, generate completions $y_1, y_2 \sim \pi_\text{ref}(\cdot \mid x)$, and annotate these pairs based on human preference to assemble an offline preference dataset $\mathcal{D} = \{x^{(i)}, y_w^{(i)}, y_l^{(i)}\}_{i=1}^N$; 2) Train the language model $\pi_\theta$ by minimizing the objective function $\mathcal{L}_\text{DPO}$, given a reference model $\pi_\text{ref}$, the dataset $\mathcal{D}$, and a chosen $\beta$ parameter.

Practically speaking, it is preferable to utilize existing, publicly released preference datasets rather than collecting new samples and corresponding human feedback. These datasets are generally constructed from samples produced by $\pi^\text{SFT}$, so we use $\pi_\text{ref} = \pi^\text{SFT}$ whenever possible. In cases where $\pi^\text{SFT}$ cannot be accessed, we instead set $\pi_\text{ref}$ to maximize the likelihood of preferred completions, namely ${\pi_\text{ref} = \argmax_{\pi}\mathbb{E}_{x, y_w \sim \mathcal{D}}\left[\log \pi(y_w \mid x)\right]}$. This initialization strategy helps address the potential discrepancy between the unavailable true reference distribution and the reference $\pi_\text{ref}$ employed in DPO. Implementation specifics and information about hyperparameter choices are described in Appendix~\ref{app:implementation}.

\section{Theoretical Analysis of DPO} This section develops additional conceptual understanding of DPO, offers formal justification, and connects the benefits of DPO with the shortcomings of actor-critic methods, like PPO~\cite{schulman2017proximal}, which are common in RLHF.

\label{sec:theory}

\subsection{Your Language Model Is Secretly a Reward Model} DPO enables joint optimization—both reward inference and policy learning—through a unified maximum likelihood objective, circumventing the need to fit an explicit reward or run a separate RL loop. Specifically, the objective in Eq.~\ref{eq:main_eq} aligns with a Bradley-Terry formulation under the reward parameterization $r^*(x, y) = \beta \log\frac{\pi^*_\theta(y \mid x)}{\pi_\text{ref}(y \mid x)}$. Optimizing the parameterized model $\pi_\theta$ in this framework is thus analogous to training a reward model, as outlined in Eq.~\ref{eq:reward_model}, following a change of variables. In the discussion below, we elaborate on the theoretical foundation for this reparameterization, demonstrate that it does not limit the family of reward models that can be learned, and prove that it enables exact recovery of the optimal policy. We start by formalizing an equivalence relation over reward functions.

\begin{definition}
Two reward functions $r(x, y)$ and $r'(x, y)$ are considered equivalent if and only if ${r(x, y) - r'(x, y) = f(x)}$ for some function $f$.
\end{definition}

It is straightforward to verify that this relation defines an equivalence, thereby dividing the set of reward functions into distinct equivalence classes. We establish the following lemmas:

\begin{lemma}\label{lemma:same_prefrence}
Within the Plackett-Luce framework, which includes the Bradley-Terry model, any pair of reward functions from the same equivalence class yield identical preference distributions.
\end{lemma}

\begin{lemma}\label{lemma:same_policy}
    Any two reward functions belonging to the same equivalence class will produce the same optimal policy when considering the constrained RL problem.
\end{lemma}

Since their proofs are direct, we refer the reader to Appendix \ref{app:lemma1} for details. The first lemma reflects a known indeterminacy present in the Plackett-Luce model family \cite{plackett1975analysis}. As a result of this indeterminacy, it is typical to introduce additional constraints to ensure identifiability when estimating the maximum likelihood parameters in Eq. \ref{eq:reward_model} \cite{bong2022generalized}. According to the second lemma, every reward function within a given equivalence class leads to the same optimal policy, so for the purpose of our objective, it is sufficient to recover any representative reward function from this class. In Appendix~\ref{app:thm1}, we prove the following theorem:

\begin{theorem}\label{thm:main}
    Assuming mild conditions, any reward class compatible with the Plackett-Luce family (including the Bradley-Terry model) admits the reparameterization ${r(x, y) = \beta \log \frac{\pi(y\mid x)}{\pi_\text{ref}(y\mid x)}}$ for some model $\pi(y\mid x)$ and a fixed reference model $\pi_\text{ref}(y \mid x)$.
\end{theorem}

\begin{sproof}
    Let $r(x, y)$ denote any reward function, with its induced optimal model $\pi_r(y \mid x)$ specified by Eq. \ref{eq:op_policy}. We demonstrate that a reward function belonging to the equivalence class of $r$ can be formulated using the above reparameterization. Define the mapping $f$ as
\begin{equation}
    f(r; \pi_\text{ref}, \beta)(x, y) = r(x, y) - \beta\log\sum_{y}\pi_\text{ref}(y\mid x)\exp\left(\frac{1}{\beta}r(x, y)\right)
\end{equation}
This operator $f$ applies a normalization to the reward using the logarithm of the partition function corresponding to $\pi_r$. Since the normalization term depends solely on the context $x$, the function $f(r; \pi_\text{ref}, \beta)(x, y)$ remains in the equivalence class of $r(x, y)$. Substituting $r$ with the right-hand side of Eq.~\ref{eq:main_eq} (valid for any reward function), it follows that $f(r; \pi_\text{ref}, \beta)(x, y) = \beta \log \frac{\pi_r(y\mid x)}{\pi_\text{ref}(y\mid x)}$. Therefore, the projection $f$ yields a member of the equivalence class of $r$ in the required parametric form, ensuring the proposed reparameterization of the reward model incurs no loss of generality.
\end{sproof}

Theorem~\ref{thm:main} can be alternatively interpreted as identifying the specific reward function chosen by the DPO reparameterization within each equivalence class. In particular, it selects the reward function that fulfills:

\begin{equation}\label{eq:lag_p}
     \sum_{y}\underbrace{\pi_\text{ref}(y\mid x)\exp\left(\frac{1}{\beta}r(x, y)\right)}_{=\pi(y\mid x)\text{, using Thm.~\ref{thm:main} reparam.}} = 1,
\end{equation}

which ensures that $\pi(y\mid x)$ is a properly normalized probability distribution, meaning all probabilities are non-negative and collectively sum to one. By comparing with Eq.~\ref{eq:op_policy}, we note that Eq.~\ref{eq:lag_p} actually defines the partition function for the optimal policy determined by the reward $r(x, y)$. The central insight behind the DPO algorithm is its capability to introduce suitable constraints into the inherently underdetermined Plackett-Luce (and especially Bradley-Terry) models of preferences. This approach maintains the expressive power of the reward model class while simultaneously making the optimal policy of Eq.~\ref{eq:op_policy} analytically manageable for any prompt $x$.

\subsection{Instability of Actor-Critic Algorithms} The proposed framework also enables analysis of the sources of instability in conventional actor-critic methods applied to RLHF, such as PPO. Adhering to the RLHF procedure, our discussion concentrates on the RL fine-tuning phase described in Section \ref{section:prelims}. This setup can be connected to the control as inference paradigm \cite{levine2018reinforcement}, as it applies to the constrained reinforcement learning formulation in Eq.~\ref{eq:RL}. Here, with a model parameterized as $\pi_{\theta}(y\mid x)$, optimization seeks to minimize $\mathbb{D}_{\text{KL}}[\pi_{\theta}(y|x) \mid \mid \pi^*(y\mid x)]$, where $\pi^*$ represents the optimal policy induced by $r_{\phi}(y, x)$ as given in Eq.~\ref{eq:optimum_model}. Some straightforward manipulations yield the following objective:

\begin{equation}\label{eq:AC}
    \max_{\pi_{\theta}}\mathbb{E}_{\pi_{\theta}(y\mid x)}\left[\underbrace{r_{\phi}(x, y) -\beta\log\sum_{y}\pi_\text{ref}(y\mid x)\exp\left(\frac{1}{\beta}r_{\phi}(x, y)\right)}_{f(r_{\phi}, \pi_\text{ref}, \beta)} - \underbrace{\beta\log\frac{\pi_{\theta}(y\mid x)}{\pi_\text{ref}(y\mid x)}}_{\text{KL}}\right]
\end{equation}

The objective considered here is identical to the one optimized in earlier studies 
\citep{ziegler2020finetuning, stiennon2022learning, bai2022training, ouyang2022training}, where a DPO-equivalent reward is used for the reward function $r_{\phi}$. In this framework, the normalization component within $f(r_{\phi}, \pi_\text{ref}, \beta)$ can be viewed as the soft value function corresponding to the reference policy $\pi_\text{ref}$. Although this term does not influence the set of optimal solutions, omitting it can lead to high-variance policy gradients, potentially causing instability during training. One way to address the normalization is by employing a learned value function; however, such an approach often introduces its own optimization challenges. Another method used in prior research involves normalizing rewards with respect to a baseline derived from human completions—this is effectively a Monte Carlo approximation of the normalization constant, based on a single sample. The DPO reparameterization, by contrast, produces a reward function that operates without needing any such baseline estimates.

\section{Experiments}
In this section, we provide an empirical assessment of DPO, investigating its effectiveness in training policies from preference data alone. We begin with a controlled experiment in text generation, designed to explore how well DPO balances reward maximization and KL-divergence minimization with respect to the reference policy, especially in comparison to widely-used preference learning techniques like PPO. We then extend our evaluation to larger model architectures and more challenging RLHF applications, such as summarization and conversational tasks. Our findings indicate that DPO typically matches or surpasses the performance of robust baselines—including RLHF with PPO and selecting the best among $N$ trajectories sampled with a trained reward model—all while requiring little to no hyperparameter tuning. A full description of our experimental methodology precedes the results, with further specifics given in Appendix~\ref{app:exp_details}.

\textbf{Tasks.} We investigate three distinct open-ended text generation tasks in our experiments. For each task, learning algorithms are trained to optimize a policy based on a dataset of preferences, denoted as $\mathcal{D}=\bigl\{x^{(i)}, y_w^{(i)}, y_l^{(i)}\bigr\}_{i=1}^N$. In the case of \textbf{controlled sentiment generation}, $x$ corresponds to the initial segment of a movie review sampled from the IMDb dataset \cite{maas-EtAl:2011:ACL-HLT2011}, and the policy is required to produce a continuation $y$ that exhibits positive sentiment. To ensure a controlled evaluation setting, we construct preference pairs over generated texts utilizing an existing sentiment classifier, ensuring $p(\text{positive}\mid x,y_w)>p(\text{positive}\mid x,y_l)$. For the SFT approach, GPT-2-large is further trained on the training partition of the IMDB dataset's reviews until performance converges (see App~\ref{app:sentiment_details} for implementation specifics). In the \textbf{summarization} task, $x$ is a Reddit forum post, and the objective is to produce a summary $y$ that captures the main arguments. In line with earlier studies, we use the Reddit TL;DR dataset \citep{volske-etal-2017-tl} and human preference annotations from \citeauthor{stiennon2022learning}. The SFT model is produced by fine-tuning on human-generated summaries of Reddit posts,\footnote{\url{https://huggingface.co/CarperAI/openai_summarize_tldr_sft}} leveraging the TRLX \citep{leandro_von_werra_2023_7790115} RLHF library. The preference data originates from the work of \citeauthor{stiennon2022learning}, using outputs from a similar, but independently trained, SFT model. Lastly, for \textbf{single-turn dialogue}, $x$ is a user prompt that may vary widely in content, ranging from scientific inquiries to personal advice requests. Here, the policy is tasked with generating a relevant and constructive reply $y$ for the user, making use of the Anthropic Helpful and Harmless dialogue dataset \citep{bai2022training}, which comprises 170k interactions between humans and an AI assistant. Each entry concludes with two alternative assistant replies, both generated by a substantial—albeit unspecified—language model, together with a label marking the response preferred by a human annotator. In this domain, because no dedicated SFT model is provided, we develop one by fine-tuning an available pre-trained language model exclusively on the preferred responses.

\textbf{Evaluation.} We employ two distinct evaluation methodologies in our experiments. For assessing how well each algorithm optimizes the constrained reward maximization task, we use a controlled sentiment generation scenario, where the performance of each algorithm is measured by plotting its tradeoff between achieved reward and KL-divergence from the reference policy. This is feasible in our setup since we have access to the true reward function (via a sentiment classifier). Conversely, in practical deployment settings where the ground-truth reward is unavailable, we instead quantify performance using the algorithm’s \textit{win rate} against a baseline policy. For this, we use GPT-4 as a stand-in for human judgment: in the summarization task, GPT-4 evaluates summary quality, and in single-turn dialogue, response helpfulness. The summarization baseline consists of reference summaries from the test data, while in dialogue, the baseline corresponds to the dataset’s preferred responses. Although recent work finds that LMs can outperform traditional automated metrics as evaluators \citep{Chen2023ExploringTU}, we further validate our approach by conducting a human study, presented in Sec.~\ref{sec:human-judgments}, to confirm the reliability of using GPT-4. Our findings indicate that GPT-4’s ratings exhibit high correlation with human judgments, often matching or surpassing the degree of agreement observed between human annotators.

\textbf{Methods.} Alongside DPO, our evaluation encompasses a range of established methods for training language models to follow human preferences. As baselines, we investigate zero-shot prompting with \textbf{GPT-J} \citep{gpt-j} for the summarization task and 2-shot prompting with \textbf{Pythia-2.8B} \citep{biderman2023pythia} for dialogue. We further compare against the \textbf{SFT} model and introduce \textbf{Preferred-FT}, a supervised fine-tuning approach using the chosen output $y_w$ either from SFT (for sentiment and summarization) or from a generic language model (for dialogue). Another supervised-like variant is \textbf{Unlikelihood}~\citep{welleck2019neural}, which trains the model to both increase the likelihood of $y_w$ and reduce that of $y_l$; an optional scaling factor $\alpha\in[0,1]$ can be applied to the unlikelihood objective. Our study also covers \textbf{PPO} \citep{schulman2017proximal}, which uses a reward model derived from preference data, and an oracle variant, \textbf{PPO-GT}, trained with direct access to the ground-truth reward signal in the controlled sentiment scenario. For PPO-GT, we utilize two forms: a standard implementation \cite{leandro_von_werra_2023_7790115} and a customized version that normalizes rewards and fine-tunes hyperparameters for superior performance; these custom adjustments are also applied to PPO with learned reward models. We additionally evaluate a \textbf{Best of $N$} strategy, where $N$ samples are drawn from either the SFT model (or Preferred-FT for dialogue), and the sample with the highest score—according to the learned reward model—is selected. Although this approach separates reward model quality from PPO optimization and achieves strong results, its computational demands are prohibitive even for moderate $N$, as every test input requires generating $N$ candidate completions.

\subsection{How well can DPO optimize the RLHF objective?}

In standard RLHF methods, the KL-constrained reward maximization framework seeks to increase the expected reward while maintaining the policy's proximity to a reference policy. Thus, evaluating algorithms requires consideration of both attained reward and the associated KL divergence; obtaining marginally better rewards at the cost of significantly increased KL divergence is generally undesirable. Figure~\ref{fig:frontier-tldr-main} illustrates the tradeoff between reward and KL for several methods within the sentiment analysis domain. For each approach, we perform several training runs, varying the key hyperparameter that controls the degree of conservativeness—specifically, the target KL $\in\{3,6,9,12\}$ for PPO, $\beta \in \{0.05,0.1,1,5\}$ for DPO, $\alpha\in\{0.05,0.1,0.5,1\}$ for unlikelihood training, and random seeds for preferred-FT—resulting in 22 runs overall. Every 100 training steps, up to convergence, we evaluate each trained policy using a suite of test prompts, calculating the mean reward according to the true reward signal and the average sequence-level KL divergence\footnote{Calculated as the aggregate of per-timestep KL values.} between the learned policy and the reference $\text{KL}\left(\pi\mid \mid \pi_\text{ref}\right)$. The results demonstrate that DPO yields a substantially more efficient tradeoff, producing policies that achieve both superior rewards and lower KL divergence compared to other methods. This outcome is noteworthy for two primary reasons: first, although both DPO and PPO target the same learning objective, DPO attains markedly better efficiency, with a reward/KL frontier that strictly outperforms PPO; second, DPO also outperforms PPO even when the latter is provided with access to ground truth rewards (PPO-GT), resulting in a more favorable frontier.

\subsection{Can DPO scale to real preference datasets?}
\label{sec:dpo-real-datasets}

We proceed by assessing how well DPO performs when fine-tuned on tasks such as summarization and single-turn dialogue. In summarization, automatic metrics like ROUGE often show weak alignment with human judgment~\citep{stiennon2022learning}, and earlier research has demonstrated that leveraging PPO for fine-tuning language models using human preferences results in more useful summaries. For evaluation, we generate model completions on the TL;DR summarization dataset’s test partition and compute the mean win rate by comparing these outputs against reference completions from the test set. Samples for each method are drawn across a range of temperature settings from 0.0 up to 1.0, with the win rates presented in Figure~\ref{fig:frontier-tldr-main} (right). DPO, PPO, and Preferred-FT are all built upon fine-tuning the same GPT-J SFT model\footnote{\url{https://huggingface.co/CarperAI/openai_summarize_tldr_sft}}. Our results indicate that DPO attains a win rate close to 61\% at temperature 0.0, outperforming PPO, which reaches around 57\% at its own optimal setting. Furthermore, DPO secures a greater maximum win rate than the strongest performance across $N$ baseline samples. It is important to mention that DPO's $\beta$ hyperparameter was not extensively optimized, implying that our reported numbers may not fully reflect DPO’s capabilities. Additionally, DPO exhibits considerably greater stability with respect to changes in sampling temperature when compared to PPO, which can regress to baseline GPT-J model performance at higher temperatures. Preferred-FT does not yield noticeable gains over the original SFT model. We also perform a direct human comparison between DPO and PPO in Section~\ref{sec:human-judgments}, showing that DPO samples generated at temperature 0.25 are favored 58\% of the time relative to PPO samples generated at temperature 0.

For single-turn dialogue evaluation, we examine the performance of various approaches on a subset of the test split from the Anthropic HH dataset \citep{bai2022training}, specifically focusing on examples containing a single exchange between a human and an assistant. In our GPT-4 based assessments, we use the test set's preferred completions as the reference to determine the win rates for each method. Since a standard SFT model tailored for this task is unavailable, our workflow begins with the pre-trained Pythia-2.8B model. We then employ Preferred-FT to train a reference model on the selected completions to ensure that generated completions remain within the model's distribution, before further training with DPO. We include a comparison to the strongest of 128 Preferred-FT completions (having observed that the Best of $N$ baseline saturates at $N=128$ in this context; refer to Appendix Figure~\ref{fig:best-of-n}), and a 2-shot prompted variant of the base Pythia-2.8B, observing that DPO meets or outperforms these baselines at optimal sampling temperatures. In addition, we evaluate an RLHF model utilizing PPO and trained on the Anthropic HH dataset\footnote{\url{https://huggingface.co/reciprocate/ppo_hh_pythia-6B}}, sourced from a widely recognized implementation\footnote{\url{https://github.com/CarperAI/trlx/tree/main/examples/hh}}, but despite our efforts, we do not identify a prompt or temperature setting that surpasses the performance of the original Pythia-2.8B. Drawing from our TL;DR findings and recognizing that both Best of 128 and PPO are trained to optimize the same reward objective, we use Best of 128 as a stand-in for PPO-level performance. In summary, DPO emerges as the sole computationally efficient technique that enhances outcomes beyond the preferred completions on the Anthropic HH dataset, matching or exceeding the performance of the far more computationally intensive Best of 128 approach. Lastly, as depicted in Figure~\ref{fig:dialogue-main}, DPO achieves its optimal performance in relatively few training steps.

\subsection{Generalization to a new input distribution}

To further assess how PPO and DPO perform under distributional shift, we test the policies trained on Reddit TL;DR summarization using a distinct dataset: the news articles found in the test split of CNN/DailyMail \citep{nallapati-etal-2016-abstractive}. For consistency, we use the optimal sampling temperatures established in the TL;DR task (0 and 0.25). Table~\ref{tab:ood} displays these outcomes. We determine the GPT-4 win rate by comparing system outputs against the human-written reference summaries using the same GPT-4 (C) prompt as before, with “forum post” replaced by “news article”. On this out-of-distribution dataset, DPO still significantly surpasses the PPO policy. These findings indicate that, despite not relying on the additional unlabeled Reddit TL;DR prompts available to PPO, DPO policies possess generalization capabilities at least comparable to those of PPO.

\subsection{Validating GPT-4 judgments with human judgments}
\label{sec:human-judgments}

A human evaluation is conducted to assess the alignment between GPT-4’s judgments and those of human raters, leveraging the TL;DR summarization results and employing two separate GPT-4 prompts. The \textbf{GPT-4 (S)} (simple) prompt asks which summary best captures the essential information in the post. Meanwhile, the \textbf{GPT-4 (C)} (concise) prompt additionally requests evaluation of conciseness; we use this because we observed that, under the simple prompt, GPT-4 tends to favor longer and more repetitive outputs than human annotators do. The full prompts are detailed in Appendix~\ref{app:prompts}. We compare three system outputs—DPO at temperature 0.25 (highest), PPO at temperature 1.0 (lowest), and SFT at temperature 0.25 (intermediate)—contrasting each with PPO at its best (greedy decoding), in order to sample a range of output qualities. In both prompt conditions, GPT-4’s agreement with human raters is on par with inter-annotator human agreement, suggesting its suitability as an evaluator when human ratings are limited (multiple human judgments are gathered for DPO vs PPO-1 only, due to constraints on human resources). Notably, the \textbf{GPT-4 (C)} prompt typically produces win rates that better correspond with human preferences; thus, it is chosen for the core results in Section~\ref{sec:dpo-real-datasets}. Further methodological details, including the rater-facing interface and the roster of human annotators, are provided in Appendix~\ref{app:human-study}.

\section{Discussion}
The paradigm of learning from preferences represents an effective and scalable method for aligning language models with desired behaviors. In this work, we have presented DPO, a straightforward approach that enables training language models using preference data without relying on reinforcement learning. DPO circumvents the need to reformulate the problem as a conventional RL task, instead leveraging a principled correspondence between reward functions and policy representations in language models. This facilitates direct alignment with human preferences through the application of a cross-entropy objective, sidestepping both reinforcement learning machinery and compromises in generality. Empirically, DPO matches or exceeds the performance of leading RLHF methods, such as those employing PPO, with minimal need for hyperparameter adjustments. As a result, DPO lowers the complexity and resource requirements associated with preference-based language model training.

\textbf{Limitations \& Future Work.} The present findings invite a range of further investigations. One key question is how the generalization capabilities of DPO-trained policies compare with approaches based on explicit reward modeling, particularly when facing out-of-distribution data. Preliminary experiments indicate that DPO achieves comparable generalization to PPO-based alternatives, but systematic evaluation is warranted. Additional topics include assessing the effectiveness of self-labeling with DPO policies in utilizing unlabeled prompt data, and examining the occurrence and effects of reward over-optimization within DPO, especially in light of the modest performance drop observed in Figure~\ref{fig:dialogue-main}-right. Furthermore, while our experiments extend to models with up to 6B parameters, scaling DPO to substantially larger, state-of-the-art models constitutes a compelling direction. Our analyses also suggest that prompt formulation impacts GPT-4-based win rate evaluations, highlighting the need for future research into optimizing automated judgment protocols. Beyond preference-driven language modeling, DPO offers the potential for broad application in training generative models across different domains.